% Source-aligned English for printed p.150.
\SourcePage{150}

\LayerTarget{tohoku:prop:2.5.2:proof:p2}
\noindent in the analogous diagram for $C'(A)$. This shows that the resulting
isomorphism is independent of the choice of $C(A)$. If we now restrict to
injective resolutions, the resulting isomorphisms are plainly natural and
commute with the connecting homomorphisms. In degree 0 they reduce to the
augmentation isomorphism, or rather its inverse, so they are indeed the
isomorphisms of Proposition 1.

\LayerTarget{tohoku:example:2.5:identity-resolution}
\noindent\textsc{Examples.} a) Let $I$ be the identity functor of $C$. A
resolving functor for $I$, also called a \textit{resolution of the identity}
in $C$, is an exact covariant functor $\mathbf C$ from $C$ to the category of
nonnegatively graded complexes in $C$, equipped with an augmentation
$A\longrightarrow C^0(A)$ that is an isomorphism from $A$ onto $Z^0C(A)$,
such that $C(A)$ is a resolution of $A$ whenever $A$ is injective. By
Proposition 2.5.1, $C(A)$ is then a resolution of $A$ for every $A$, since
$I$ is exact and $R^nI=0$ for $n>0$.

Let $\mathbf C$ be a resolution of the identity in $C$, and let $F$ be a
covariant left exact functor from $C$ to $C'$. Suppose
$R^nF(C^i(A))=0$ for every $n>0$, so the $C^i(A)$ are $F$-acyclic.
\textit{Then $\mathbf F(A)=F(C(A))$ is a resolving functor for $F$.} Indeed,
it is exact: from an exact sequence
$0\longrightarrow A'\longrightarrow A\longrightarrow A''\longrightarrow0$
one obtains
$0\longrightarrow C^i(A')\longrightarrow C^i(A)\longrightarrow
C^i(A'')\longrightarrow0$, and since $R^1F(C^i(A'))=0$, also
\[
0\longrightarrow F(C^i(A'))\longrightarrow F(C^i(A))
\longrightarrow F(C^i(A''))\longrightarrow0.
\]
This proves (i). Condition (ii) follows because $F$ is left exact and
$\mathbf C$ is exact. Condition (iii) follows because, with $C(A)$ an
$F$-acyclic resolution of $A$, one has $H^n(F(C(A)))=R^nF(A)$, which vanishes
for $n>0$ when $A$ is injective.

A convenient way to construct a resolution of the identity in $C$ is to take
an exact functor $C^0$ from $C$ to $C$ and a natural monomorphism
$A\longrightarrow C^0(A)$. Define the $C^n(A)$ recursively, together with
$d^{n-1}:C^{n-1}\longrightarrow C^n$, by
\[
C^1(A)=C^0(C^0(A)/\operatorname{Im}A)
\]
and, for $n\geq2$,
\[
C^n(A)=C^0(C^{n-1}(A)/\operatorname{Im}C^{n-2}(A)).
\]
Here $d^{n-1}$ is defined using the augmentation homomorphism from
$Q=C^{n-1}(A)/\operatorname{Im}C^{n-2}(A)$ to $C^0(Q)$. This gives a
resolution of the identity. For the $C^i(A)$ to be $F$-acyclic, it suffices
that the $C^0(A)$ be so.

\LayerTarget{tohoku:example:2.5:hom-projective-resolution}
\noindent b) Let $P=(P_n)$ be a projective resolution of an object $A$ of
$C$, and let $F(B)=\operatorname{Hom}(A,B)$ be the functor from $C$ to
abelian groups. Then $\operatorname{Hom}(P,B)$ is a resolving functor for
$F$.

\LayerTarget{tohoku:construction:resolving-functor-spectral-sequences:p1}
We now show how the most important spectral sequences can be computed by
means of resolving functors. Let $\mathbf F$ be a resolving functor for a
left exact functor $F:C\longrightarrow C'$. Let $K$ be a complex in $C$
bounded on the left, and consider the bicomplex
$\mathbf F(K)=(\mathbf F^q(K^p))_{p,q}$. It depends naturally on $K$, as do
its two spectral sequences and their abutment. By the exactness of
$\mathbf F$ and Proposition 2.5.1, the initial terms are
\LayerTarget{tohoku:eq:2.5.3}
\[
\tag{2.5.3}
{}_{\mathrm I}E_2^{p,q}(\mathbf F(K))=H^p(R^qF(K)),
\qquad
{}_{\mathrm{II}}E_2^{p,q}(\mathbf F(K))=R^pF(H^q(K)).
\]
These are natural isomorphisms with the initial terms of the

\TohokuContinuation{151}
